\section{Namensdienst und Worker-Management}
\label{sec:namensdienst}

\subsection{Service Discovery -- Konzept}

In TaskGrid+ werden Worker-Adressen \textbf{nirgends statisch} konfiguriert.
Der Dispatcher kennt zur Build-Zeit nur den Namensdienst.
Welche Worker für welchen Task-Typ verfügbar sind, erfährt der Dispatcher
ausschließlich durch Abfragen an den Namensdienst (\texttt{LookupWorker}).

Dieses Prinzip ist in der Aufgabenstellung (§4.2) explizit gefordert
und wurde in allen Dispatcher-Konfigurationsdateien (Dockerfile,
docker-compose.yml) durchgesetzt:
Worker-Adressen erscheinen in keiner Konfigurationsdatei des Dispatchers.

\subsection{Erwartetes Worker-Verhalten (Schnittstelle)}

Damit der Dispatcher korrekt funktioniert, müssen Worker folgendes
Protokoll einhalten -- dies ist die Schnittstelle, auf die sich der
Dispatcher verlässt:

\begin{itemize}
  \item \textbf{Registrierung beim Start}: Jeder Worker ruft
        \texttt{RegisterWorker} am Namensdienst auf mit seiner
        \texttt{worker\_id}, den unterstützten \texttt{task\_types},
        seiner \texttt{address} (Docker-Service-Name) und seinem \texttt{port}.
  \item \textbf{Regelmäßige Heartbeats}: Worker senden periodisch
        \texttt{SendHeartbeat}-Nachrichten, damit der Namensdienst
        sie als aktiv führt.
  \item \textbf{Deregistrierung beim Beenden}: Worker rufen
        \texttt{DeregisterWorker} auf, damit sie sofort aus der Registry
        entfernt werden (ohne auf das Heartbeat-Timeout zu warten).
  \item \textbf{Ergebnis-Rückmeldung}: Nach Verarbeitung eines Tasks
        ruft der Worker \texttt{ReturnResult} am Dispatcher auf
        (nicht am Namensdienst).
\end{itemize}

Die Implementierung dieses Worker-Verhaltens liegt außerhalb des
in dieser Dokumentation beschriebenen Aufgabenbereichs.

\subsection{Dispatcher-seitiger Namensdienst-Client (Issue \#14)}

Der \texttt{NamensdienstClient} (\texttt{src/dispatcher/namensdienst\_client.py})
kapselt alle gRPC-Aufrufe des Dispatchers an den Namensdienst.
Die Verbindungsadresse wird ausschließlich aus Umgebungsvariablen gelesen:

\begin{lstlisting}[language=Python,caption={NamensdienstClient -- Konfiguration aus Umgebungsvariablen}]
host = os.environ.get("NAMENSDIENST_HOST", "namensdienst")
port = os.environ.get("NAMENSDIENST_PORT", "50052")
self._target = f"{host}:{port}"
self._channel = grpc.insecure_channel(self._target)
self._stub = taskgrid_pb2_grpc.NamingServiceStub(self._channel)
\end{lstlisting}

Die einzige nach außen exponierte Methode ist
\texttt{lookup\_worker(task\_type, request\_id)}.
Sie ist \textbf{defensiv} implementiert:
Bei einem gRPC-Fehler (Namensdienst nicht erreichbar, Timeout)
wird keine Exception geworfen, sondern eine leere Liste zurückgegeben.

\begin{lstlisting}[language=Python,caption={Fehlertolerante LookupWorker-Implementierung}]
try:
    response = self._stub.LookupWorker(
        taskgrid_pb2.LookupRequest(task_type=task_type),
        timeout=5.0,
    )
except grpc.RpcError as e:
    log_event(logger, "error", "LOOKUP_WORKER_failed",
              request_id=request_id,
              task_type=task_type,
              error=str(e.code()))
    return []   # kein Absturz -- Dispatcher wartet und versucht erneut
\end{lstlisting}

Der Dispatcher behandelt eine leere Liste wie \glqq{}kein Worker verfügbar\grqq{}
und re-enqueued den Task nach \texttt{DISPATCH\_NO\_WORKER\_RETRY\_SECONDS}
(default: 2\,s).
Strukturiertes Logging gemäß §13 wird in beiden Fällen (Erfolg und Fehler)
erzeugt:

\begin{lstlisting}[language={},caption={Beispiel-Logs des NamensdienstClients}]
[dispatcher.namensdienst_client] task_type=sum event=LOOKUP_WORKER_success
                                  worker_count=3
[dispatcher.namensdienst_client] task_type=sum event=LOOKUP_WORKER_failed
                                  error=StatusCode.UNAVAILABLE
[dispatcher.namensdienst_client] task_type=sum event=LOOKUP_WORKER_no_workers
\end{lstlisting}

\subsection{Worker-Auswahl per Round-Robin (Issue \#14)}

Wenn \texttt{LookupWorker} mehrere Worker zurückgibt, wählt der
\texttt{RoundRobinSelector} (\texttt{src/dispatcher/worker\_selector.py})
einen davon aus.
Pro Task-Typ wird ein eigener Zähler geführt, sodass zwei Typen
sich nicht gegenseitig beeinflussen:

\begin{lstlisting}[language=Python,caption={Round-Robin-Selektion (worker\_selector.py)}]
class RoundRobinSelector:
    def select(self, task_type: str,
               workers: list[WorkerInfo]) -> Optional[WorkerInfo]:
        if not workers:
            return None
        with self._lock:
            idx = self._counters[task_type] % len(workers)
            self._counters[task_type] += 1
        return workers[idx]
\end{lstlisting}

Der Counter ist thread-sicher (\texttt{threading.Lock}) und wird bei
jedem Aufruf inkrementiert.
Bei drei Workern verteilen sich sechs aufeinanderfolgende Tasks
gleichmäßig: Worker\,0, Worker\,1, Worker\,2, Worker\,0, Worker\,1, Worker\,2.

\subsection{Verhalten bei fehlendem Worker}

Der Dispatch-Loop in \texttt{src/dispatcher/dispatch\_loop.py}
behandelt den Fall, dass kein Worker verfügbar ist, explizit:

\begin{lstlisting}[language=Python,caption={Kein-Worker-Behandlung im Dispatch-Loop}]
worker = self._selector.select(task.task_type, workers)

if worker is None:
    log_event(logger, "warning", "DISPATCH_no_worker",
              task_id=task.task_id, task_type=task.task_type)
    time.sleep(_NO_WORKER_RETRY_SECS)
    self._queue.enqueue(task)   # Task wird erneut eingereiht
    return
\end{lstlisting}

Der Task verbleibt im Zustand \texttt{QUEUED}, sein \texttt{retry\_count}
wird \emph{nicht} inkrementiert.
Dies ist der einzige Pfad durch das System, bei dem ein Task
\texttt{QUEUED} bleibt ohne einen Zustandswechsel auszulösen.

\subsection{Skalierung aus Dispatcher-Sicht}

Werden mehrere Instanzen desselben Worker-Typs gestartet
(z.\,B. \texttt{docker compose up --scale worker-sum=3}),
registrieren sich alle drei beim Namensdienst mit
unterschiedlichen \texttt{worker\_id}-Werten aber demselben
Service-Namen (\texttt{worker-sum}).

Aus Dispatcher-Sicht ändert sich nichts am Code:
\texttt{LookupWorker} gibt drei Einträge zurück, und der
Round-Robin-Selector verteilt Tasks automatisch auf alle drei.
Der Dispatcher selbst wird nicht angepasst.
